%% ma: L10-B15-0105
%% lop: 10
%% chuong: 6
%% bai: 15
%% dang: 10.15.4
%% loai: DS
%% muc: [NB, TH, NB, NB]
%% dap_an: "SĐĐS"
%% nguon: "Ảnh giáo viên gửi 25/11/2026 (lô 2): Bài tập trắc nghiệm – Toán 10 Bài 15 Hàm số (sự biến thiên, đồ thị), Phần II Câu 1"
%% kien_thuc: "Lớp 10 Bài 15 (đồ thị, đồng biến nghịch biến); hàm số y=ax^2, phương trình bậc hai (lớp 9)"
%% trang_thai: chinh-thuc
%% ghi_chu: "Ý c gốc ghi 'Đồ thị của hàm số ... nghịch biến', đã sửa lời thành 'Hàm số ... nghịch biến'."
\begin{ex}
Cho đồ thị các hàm số $y=-2x+3$; $y=2x^2$. Xét tính đúng sai của các khẳng định sau:
\begin{center}\begin{tikzpicture}[scale=.8]
\draw[truc] (-2,0)--(2,0) node[below]{$x$};
\draw[truc] (0,-.6)--(0,6) node[left]{$y$};
\node[below] at (-1,0) {\footnotesize$-1$}; \node[below right] at (0,0) {\footnotesize$O$};
\draw[thick,domain=-1.7:1.7,smooth] plot (\x,{2*(\x)^2}) node[right] {\footnotesize$y=2x^2$};
\draw[thick,domain=-1.4:1.45] plot (\x,{-2*(\x)+3});
\node[right] at (0.95,1.1) {\footnotesize$y=-2x+3$};
\end{tikzpicture}\end{center}
\choiceTF
{Đồ thị hàm số $y=-2x+3$ là một đường cong.}
{\True Đồ thị hàm số $y=-2x+3$ cắt đồ thị hàm số $y=2x^2$ tại hai điểm.}
{\True Hàm số $y=-2x+3$ nghịch biến trên $\mathbb R$.}
{Hàm số $y=2x^2$ nghịch biến trên khoảng $(0;+\infty)$.}
\loigiai{a) Đồ thị là đường thẳng: Sai. b) $2x^2=-2x+3\Leftrightarrow2x^2+2x-3=0$ có $\Delta'=7>0$ nên có hai giao điểm: Đúng. c) $a=-2<0$: Đúng. d) Với $0<x_1<x_2$ thì $2x_1^2<2x_2^2$ nên hàm số đồng biến: Sai.}
\end{ex}
